\scanpage{40}
Enfin, supposons maintenant que $K$ soit isomorphe
$\mathbb R$ [on peut supposer $K=\mathbb R$, car $\mathbb R$ n'a pas d'
autom. non trivial] soit $K'$ une clôture
algébrique de $K$, $\pi=\mathbb Z/2\mathbb Z$ son groupe de Galois,
$X$ lisse \deleted{\unsure{algébrique|projectif}} [projectif] sur $K$. Alors on peut appliquer
ce qui précède à $X_{K'}$ sur $K'$,
\deleted{$H^*(X_{K'\,\mathrm{top}},\mathbb R)\unsure{\xrightarrow{\sim}|\longrightarrow} H^*(X)\otimes_K K'$}, on trouve
\[
\left[\varphi_{K'}:\coprod_{p,q}H^q(X,\Omega_X^p)\otimes_K K'
 \xrightarrow{\sim}H^*(X_{\deleted{\unsure{K'}}})\otimes_K K'\right]
\]
et d'après ce qui a été vu plus haut,
\unsure{cet|ces} isomorphisme\deleted{\unsure{s}} \uline{commute\deleted{\unsure{nt}} à} $\pi$, donc provient
d'un isomorphisme
\[
 \varphi_1:\coprod_{p,q} H^q(X,\Omega_X^p)\xrightarrow{\sim}H^*(X)
\]

% margin note written sideways
\margin{On a aussi un $\mathbb R$-isom canonique $(\pi=\mathbb Z/2\mathbb Z)$
$H^*(X_{\mathbb C\,\mathrm{top}},\mathbb R)\otimes_{\mathbb R}\mathbb C\xrightarrow{\sim}H^*(X)\otimes_{\mathbb R}\mathbb C$,
où $\pi$ opère à gauche via son opération sur $X_{\mathbb C\,\mathrm{top}}$,
\uline{et} sur $\mathbb C$ dans $\otimes_{\mathbb R}\mathbb C$
d'où en prenant les invariants sous $\pi$ une décomposition
canonique $H^*(X)\simeq H^*(X)^+\oplus H^*(X)^-$,
\deleted{$\unsure{H^*(X,\mathbb R)}$}
$H^*(X)^+\simeq H^*(X_{\mathbb C\,\mathrm{top}},\mathbb R)^\pi\simeq H^*(X_{\mathrm{top}}=X_{\mathbb C\,\mathrm{top}}/\pi_0\,;\mathbb R)$,
$H^*(X)^-=H^*(X_{\mathbb C\,\mathrm{top}},\mathbb R)^-$.}

En résumé, chaque fois que $X$ est
projectif non singulier sur un corps
\deleted{\unsure{top. ou|top.}} top. loc. \unsure{cpt|cpt.} complet archimédien,
on a une décomposition $\varphi$ comme
dessus de $H^*(X)$, compatible avec la
filtration. \deleted{Si} On pose
\[
H^{p,q}(X)=\varphi\big(H^q(X,\Omega_X^p)\big)\subset \deleted{\unsure{F|E}}F^p H^q(X)
\]

Si $K$ n'est pas topologique, alors pour
tout choix d'une topologie \unsure{forterie|forte|localement} archimédienne
$T$ sur $K$, on déduit sur
% archive sheet number: 39
